\begin{examplebox}
	\textbf{\textcolor{CExample}{例 1（基础）}}

	\textbf{\textcolor{CExample}{题目：}}
	设正数 $a,b$ 满足
	\[
		\frac{a-b}{\ln a-\ln b}=3,
	\]
	求证 $a+b>6$ 且 $ab<9$。

	\textbf{\textcolor{CExample}{解题步骤：}}
	\begin{description}[style=nextline, leftmargin=2.8em, labelsep=0.5em, itemsep=0.45em]
		\item[\textbf{第一步（确认条件）：}]
		      由题设，$a,b>0$。由于分母 $\ln a-\ln b$ 不能为 $0$，所以 $a\neq b$。等式左端即为对数均值形式，常数 $c=3$。
		      \par\textbf{理由：} 此步确认满足定理的使用前提。

		\item[\textbf{第二步（套和结论）：}]
		      由定理 $x_1+x_2>2c$，代入 $a,b$ 得
		      \[
			      a+b>2\times3=6.
		      \]
		      \par\textbf{理由：} 直接套用对数均值结论。

		\item[\textbf{第三步（套积结论）：}]
		      同理，由 $x_1x_2<c^2$ 得
		      \[
			      ab<3^2=9.
		      \]
		      \par\textbf{理由：} 直接套用对数均值结论。
	\end{description}

	\textbf{\textcolor{CExample}{关键结论：}}
	\textbf{$a+b>6$，$ab<9$。}

	\medskip
	\textbf{\textcolor{CExample}{例 2（稍变形）}}

	\textbf{\textcolor{CExample}{题目：}}
	已知函数 $f(x)=\ln x-x$，正实数 $x_1,x_2$ 满足 $f(x_1)=f(x_2)$ 且 $x_1\neq x_2$，求证
	\[
		x_1+x_2>2,
		\qquad
		x_1x_2<1.
	\]

	\textbf{\textcolor{CExample}{解题步骤：}}
	\begin{description}[style=nextline, leftmargin=2.8em, labelsep=0.5em, itemsep=0.45em]
		\item[\textbf{第一步（化标准形式）：}]
		      由 $f(x_1)=f(x_2)$ 得
		      \[
			      \ln x_1-x_1=\ln x_2-x_2.
		      \]
		      移项得
		      \[
			      \ln x_1-\ln x_2=x_1-x_2,
		      \]
		      因此
		      \[
			      \frac{x_1-x_2}{\ln x_1-\ln x_2}=1.
		      \]
		      即 $c=1$。
		      \par\textbf{理由：} 将给定条件变形为结论所需结构。

		\item[\textbf{第二步（检查条件）：}]
		      $x_1,x_2>0$，且 $x_1\neq x_2$，符合定理条件。
		      \par\textbf{理由：} 对数有意义且分母不为零。

		\item[\textbf{第三步（套用结论）：}]
		      直接得
		      \[
			      x_1+x_2>2c=2,
			      \qquad
			      x_1x_2<c^2=1.
		      \]
		      \par\textbf{理由：} 由定理直接推出。
	\end{description}

	\textbf{\textcolor{CExample}{关键结论：}}
	\textbf{$x_1+x_2>2$，$x_1x_2<1$。}

	\medskip
	\textbf{\textcolor{CExample}{例 3（综合应用）}}

	\textbf{\textcolor{CExample}{题目：}}
	设 $a>0$，函数 $f(x)=\ln x-ax$ 有两个不同的零点 $x_1,x_2$，求证
	\[
		x_1+x_2>\frac{2}{a},
		\qquad
		x_1x_2<\frac{1}{a^2}.
	\]

	\textbf{\textcolor{CExample}{解题步骤：}}
	\begin{description}[style=nextline, leftmargin=2.8em, labelsep=0.5em, itemsep=0.45em]
		\item[\textbf{第一步（零点转化条件）：}]
		      由零点定义可知
		      \[
			      f(x_1)=0,
			      \qquad
			      f(x_2)=0.
		      \]
		      即
		      \[
			      \ln x_1=ax_1,
			      \qquad
			      \ln x_2=ax_2.
		      \]
		      两式相减得
		      \[
			      \ln x_1-\ln x_2=a(x_1-x_2),
		      \]
		      故
		      \[
			      \frac{x_1-x_2}{\ln x_1-\ln x_2}=\frac{1}{a}.
		      \]
		      即 $c=\dfrac{1}{a}$。
		      \par\textbf{理由：} 由零点等式导出标准形式，且 $a>0$，所以 $c>0$。

		\item[\textbf{第二步（确认条件）：}]
		      由零点存在性，$x_1,x_2>0$，且两个零点不同，所以 $x_1\neq x_2$，符合定理前提。
		      \par\textbf{理由：} 保证 $\ln x$ 有意义，且对数均值分母不为零。

		\item[\textbf{第三步（套用结论）：}]
		      利用 $x_1+x_2>2c$ 得
		      \[
			      x_1+x_2>\frac{2}{a}.
		      \]
		      利用 $x_1x_2<c^2$ 得
		      \[
			      x_1x_2<\frac{1}{a^2}.
		      \]
		      \par\textbf{理由：} 定理的直接应用。
	\end{description}

	\textbf{\textcolor{CExample}{关键结论：}}
	\textbf{$x_1+x_2>\dfrac{2}{a}$，$x_1x_2<\dfrac{1}{a^2}$。}
\end{examplebox}